% =========================================================
% SECCIÓN 3.2 - FUNDAMENTOS DEL MODELO DEB
% =========================================================

El Dynamic Energy Budget (DEB) constituye un marco teórico unificador para describir el metabolismo de organismos heterótrofos mediante balances de masa y energía a lo largo de su ciclo de vida \parencite{eriksenDynamicModellingFeed2022}. A diferencia de modelos empíricos de crecimiento ---que ajustan curvas logísticas o de Gompertz sin desagregar los flujos subyacentes---, el DEB explicita cómo la materia ingerida se distribuye entre mantenimiento, crecimiento somático, reserva reproductiva y disipación metabólica, expresando cada flujo en unidades de carbono por unidad de tiempo.

En el contexto de la bioconversión con \emph{Hermetia illucens}, el modelo DEB adquiere una función predictiva dual: no solo estima la acumulación de biomasa larval, sino que cuantifica la tasa de respiración ---y por ende la emisión de \ch{CO2}--- como resultado inevitable de los costos de mantenimiento y síntesis. Esta sección resume los principios del DEB escalado a larvas de BSF y su adaptación al reactor batch descrito en el Capítulo~\ref{chap:instrumentacion}.

% ---------------------------------------------------------
\subsection{Estructura de compartimentos de carbono}
\label{subsec:cap3_compartimentos}
% ---------------------------------------------------------

Siguiendo la parametrización de \textcite{eriksenDynamicModellingFeed2022}, se define un sistema de cinco compartimentos de carbono interconectados. Las variables $E$, $B$ y $L$ se expresan primero como contenidos específicos por larva (\si{\gram C \per larva}) y luego se escalan al reactor completo; $S$ y $C_{\ch{CO2}}$ son variables extensivas del reactor (\si{\gram C}):

\begin{enumerate}
    \item \textbf{Sustrato orgánico disponible} ($S$): masa total de carbono orgánico en el reactor susceptible de ser ingerida por las larvas, expresada en \si{\gram C}. Incluye los polímeros degradables (carbohidratos, proteínas, lípidos) pero excluye la fracción lignocelulósica recalcitrante, que en el sustrato de pulpa de café y cáscaras de frutas (Sección~\ref{sec:cap2_biomassa}) representa menos del $15\%$ de la masa total \parencite{brunoValorizationOrganicWaste2025}.

    \item \textbf{Reserva de asimilados} ($E$): carbono absorbido por el intestino larval y almacenado temporalmente antes de su destino a metabolismo energético o biosíntesis. Eriksen la identifica como el ``banco'' de carbono que amortigua la variabilidad entre ingesta y demanda metabólica. Su unidad base es \si{\gram C \per larva}; el contenido total del reactor es $E_{\mathrm{tot}} = N_{\mathrm{larvas}} \cdot E$.

    \item \textbf{Biomasa estructural} ($B$): carbono incorporado en tejidos somáticos (músculo, cutícula, sistema nervioso, tracto digestivo). Esta variable se correlaciona directamente con la biomasa húmeda medida destructivamente cada 48~h en el protocolo del Capítulo~\ref{chap:instrumentacion}, mediante un factor de conversión $\rho_{\mathrm{conv}}$ (\si{\gram C \per \gram\;biomasa\;húmeda}) que se calibrará localmente. Unidad base: \si{\gram C \per larva}.

    \item \textbf{Lípidos de reserva} ($L$): fracción de carbono destinada al almacenamiento energético de largo plazo, predominante en estadios tardíos (L5--prepupa). \textcite{grausaMetabolicModelingHermetia2023} documentan que esta fracción puede alcanzar hasta el $35\%$ de la biomasa seca en sustratos ricos en carbohidratos, y constituye un sumidero de carbono paralelo a la estructura. Unidad base: \si{\gram C \per larva}.

    \item \textbf{Dióxido de carbono acumulado} ($C_{\ch{CO2}}$): carbono total emitido como \ch{CO2} por respiración larval y microbiana, expresado en \si{\gram C}. En el reactor batch cerrado, esta variable se vincula directamente con la concentración medida por el sensor NDIR mediante la ecuación de conversión derivada en la Sección~\ref{sec:cap2_conversion_flujo}.
\end{enumerate}

La elección del carbono como \emph{currency} metabólica ---en lugar de la energía en joules--- responde a dos criterios prácticos. Primero, los sensores NDIR miden directamente la concentración de \ch{CO2}, que es un \emph{proxy} lineal del carbono respirado si se desprecia la fracción disuelta en el sustrato (menor al $2\%$ en condiciones de pH~5.5--6.5, según el equilibrio \ch{CO2}/\ch{HCO3-} a $30^{\circ}$C). Segundo, los balances de masa de carbono permiten validación cruzada con el método de \emph{mass balance} estándar en estudios de BSF \parencite{parodiBioconversionEfficienciesGreenhouse2020,rossiEstimatingDynamicsGreenhouse2024a}.

% ---------------------------------------------------------
\subsection{Flujos metabólicos y producción de \texorpdfstring{\ch{CO2}}{CO2}}
\label{subsec:cap3_flujos}
% ---------------------------------------------------------

El metabolismo larval se descompone en cinco flujos principales. Salvo indicación explícita, las tasas $\dot{I}, \dot{A}, \dot{M}, \dot{G}_B, \dot{G}_L$ y $\dot{D}$ se expresan como flujos específicos por larva (\si{\gram C \per larva \per \day}), de modo que el balance total del reactor se obtiene multiplicando por $N_{\mathrm{larvas}}$:

\begin{description}
    \item[Ingesta ($\dot{I}$)] Velocidad a la que el carbono del sustrato pasa al tracto digestivo. Está limitada por la densidad larval, la humedad del sustrato y la disponibilidad física de alimento. \textcite{dienerConversionOrganicMaterial2011} reportan tasas máximas de 100~mg alimento/larva/día en condiciones óptimas.

    \item[Asimilación ($\dot{A}$)] Fracción de la ingesta que cruza la pared intestinal y se incorpora a la reserva $E$. La eficiencia de asimilación $\kappa_A$ depende de la calidad del sustrato; valores entre $0.55$ y $0.80$ son típicos para residuos agroindustriales \parencite{guillaumeAsymptoticEstimatedDigestibility2023}.

    \item[Mantenimiento somático ($\dot{M}$)] Costo basal de preservar la viabilidad celular, la homeostasis iónica y la motilidad. Es proporcional a la biomasa estructural $B$ y modulado por la temperatura según una cinética tipo Arrhenius o Sharpe--Schoolfield (Sección~\ref{sec:cap3_cinetica}). Eriksen estima que el mantenimiento consume entre el $15\%$ y el $25\%$ del carbono asimilado en larvas de BSF.

    \item[Síntesis de biomasa ($\dot{G}_B$) y de lípidos ($\dot{G}_L$)] Flujos de crecimiento que translocan carbono desde la reserva $E$ hacia los compartimentos estructural y lipídico. La prioridad entre $\dot{G}_B$ y $\dot{G}_L$ cambia con el estadio larval: domina la síntesis de estructura en L3--L4, mientras que la lipogénesis se acelera en L5 previo a la metamorfosis \parencite{eriksenDynamicModellingFeed2022,grausaMetabolicModelingHermetia2023}.

    \item[Disipación ($\dot{D}$)] Carbono perdido irrevocablemente como \ch{CO2} durante los procesos de mantenimiento y síntesis. En la parametrización adaptada de Eriksen, la respiración específica total es la suma de tres términos disipativos:
    \begin{equation}
        \dot{D}_{\mathrm{ind}}(t) = \dot{M}(t) + \bigl(1-\kappa_G\bigr)\dot{G}_B(t) + \bigl(1-\kappa_L\bigr)\dot{G}_L(t),
        \label{eq:cap3_disipacion_total}
    \end{equation}
    donde $\kappa_G$ y $\kappa_L$ representan las eficiencias de conversión de reserva en biomasa estructural y lípidos, respectivamente. El término $(1-\kappa_G)$ captura el costo de síntesis (\emph{overhead}), que en insectos holometábolos suele situarse entre $0.2$ y $0.4$ \parencite{nielsenMetabolicPerformanceMealworms2025}.
\end{description}

La producción de \ch{CO2} medida por los sensores NDIR ---expresada como flujo másico $\dot{m}_{\ch{CO2}}(t)$ en \si{\gram\per\minute}--- se identifica con la disipación total del reactor (larvas + microbiota), escalada mediante el factor molar carbono--\ch{CO2}:
\begin{equation}
    \dot{m}_{\ch{CO2}}(t) = \frac{M_{\ch{CO2}}}{M_{\ch{C}}}\,\Bigl[\,N_{\mathrm{larvas}}\,\dot{D}_{\mathrm{ind}}(t) + \dot{D}_{\mu}(t)\Bigr],
    \label{eq:cap3_puente_sensor}
\end{equation}
donde $M_{\ch{CO2}} = 44.01$~\si{\gram\per\mol} y $M_{\ch{C}} = 12.01$~\si{\gram\per\mol} son las masas molares del dióxido de carbono y del carbono elemental, respectivamente. El número de larvas $N_{\mathrm{larvas}}$ se deriva de la biomasa inicial del protocolo experimental ($\approx 700$~larvas por reactor, Sección~\ref{sec:cap2_biomassa}) y se considera constante durante el ciclo dado que no se registró mortalidad significativa en las réplicas. El término $\dot{D}_{\mu}(t)$ representa la respiración microbiana del sustrato no ingerido; su cinética se modelará en la Sección~\ref{sec:cap3_formulacion_ode} como fracción del sustrato remanente $S(t)$.

La ecuación~\eqref{eq:cap3_puente_sensor} constituye el puente cuantitativo entre la teoría bioenergética y el dato crudo del Capítulo~\ref{chap:instrumentacion}: permite interpretar una pendiente de concentración de 200~ppm en 30~min no como un número aislado, sino como la manifestación física de los costos metabólicos de un conjunto larval en un estado fisiológico determinado.

Es relevante notar que el modelo DEB, en su formulación estándar para BSF, asume metabolismo aeróbico exclusivo. Por tanto, la generación de \ch{CH4} no aparece en los balances de carbono de las EDO; su presencia transitoria se trata como una perturbación estocástica que indica desviación del régimen aeróbico, tema que se aborda en el Capítulo~\ref{chap:metricas} mediante el clasificador de eventos de anaerobiosis.

En síntesis, los compartimentos de carbono y los flujos metabólicos descritos en esta sección proporcionan la arquitectura conceptual sobre la cual se construye el sistema de ecuaciones diferenciales de la Sección~\ref{sec:cap3_formulacion_ode}. La consistencia del balance de carbono ---que se verifica formalmente en la Sección~\ref{sec:cap3_analisis}--- garantiza que toda unidad de carbono ingerida se \emph{accountabilice} como biomasa, lípidos, \ch{CO2} o sustrato residual, cerrando el ciclo de masa requerido para la validación del modelo contra datos experimentales.